\documentclass[11pt,a4paper]{moderncv}
\moderncvstyle{banking} 
\moderncvcolor{blue}

% Aumentamos el 'scale' a 0.82 para solucionar el "Overfull \vbox" y aprovechar mejor la página
\usepackage[scale=0.82]{geometry}
% Reintegramos hyperref para que reconozca \hypersetup
\usepackage{hyperref}

\hypersetup{
	colorlinks=true,
	urlcolor=color1,
	linkcolor=color1
}

\name{Matihus}{Molina Larios}
\address{Lima, Perú}{}
\email{matihus.molina@unmsm.edu.pe}
\social[linkedin]{maticus}
\social[github]{Maticus-7}
\phone{+51 903 003 327}

\begin{document} 
	\makecvtitle
	
	\section{Perfil profesional}
	\cvitem{}{Estudiante de último año de Ciencias Matemáticas en la UNMSM con un sólido rigor analítico enfocado en Inteligencia Artificial y Ciencia de Datos. Cuento con sólida formación matemática y estadística aplicada en el procesamiento de datos estructurados y el desarrollo de modelos predictivos de Machine Learning. Actualmente fortalezco mis habilidades en un programa de Advanced Data Science. Busco mi primera experiencia pre-profesional para aportar valor técnico en la validación, optimización y desarrollo de modelos de IA, transformando datos complejos en decisiones estratégicas en el sector financiero.}
	
	\section{Habilidades Técnicas y Tecnologías}
	\cvitem{Lenguajes}{Python (Pandas, NumPy, Matplotlib, Seaborn), SQL, R, MATLAB, LEAN4.}
	\cvitem{Machine Learning}{Scikit-learn, PyTorch, Modelado Predictivo y Prescriptivo, Validación de Modelos, Procesamiento de Datos Estructurados, Optimización Matemática, Métodos Numéricos.}
	\cvitem{Herramientas}{Git, GitHub, Jupyter Notebooks, Power BI, Microsoft Excel (Intermedio), \LaTeX.}
	
	\section{Educación}
	\cventry{03/2022 -- 12/2026 (estimado)}{\textnormal{Estudiante de Bachillerato en Ciencias Matemáticas}}{Universidad Nacional Mayor de San Marcos}{\textnormal{Lima, Perú}}{}{(Perteneciente al décimo superior)}
	\cventry{2018 -- 2021}{\textnormal{Egresado del Programa de Inglés Integral}}{Centro Cultural LIDEMPERU}{\textnormal{Lima, Perú}}{}{}
	
	\section{Experiencia en Proyectos e Investigación}
	\cventry{02/2025 -- 03/2025}{Predicción de Depresión en Estudiantes}{Curso: Matemáticas para ML y DS (UNMSM)}{}{}{
		\begin{itemize}
			\item Desarrollé un flujo de trabajo de Machine Learning utilizando Python y Scikit-learn para predecir niveles de depresión basados en variables socioeconómicas y académicas.
			\item Implementé el procesamiento integral de datos estructurados: imputación de nulos, codificación de variables categóricas (One-Hot Encoding) y normalización.
			\item Entrené y validé modelos supervisados comparando métricas clave como Accuracy, Precision, Recall, F1-Score y análisis de matriz de confusión para optimizar la toma de decisiones.
			\item \href{https://github.com/Maticus-7/Mental-Healthy/tree/2bc0dac3f89d9187dfa37b6cd963d69256551167/Depression/University}{Haga click aquí para ver el repositorio en GitHub}
	\end{itemize}}
%Dar un fin:
	\cventry{01/2025 -- Presente}{Portafolio: Aprendizaje Estadístico y Optimización Aplicada}{Proyectos de Especialización (Futura \& SaC$^3$)}{}{}{
		\begin{itemize}
			\item Desarrollé un repositorio enfocado en la resolución de problemas de probabilidad, estadística computacional y optimización utilizando Python. %Combinar con el tercero
			\item Implementé y documenté algoritmos prácticos de optimización numérica analizados en el SaC$^3$ y flujos de análisis de datos basados en el programa Advanced Data Science de Futura.
			\item Apliqué buenas prácticas de control de versiones con Git y GitHub para la estructuración de código limpio y reproducible. %tercero xd
			\item \href{https://github.com/Maticus-7/Data-Science}{Haga click aquí para ver el repositorio en GitHub}
	\end{itemize}}
	
	\cventry{10/2025 -- 12/2025}{Numerical solution for differential equations of Lane-Emden type}{Modelamiento de Cuerpos Estelares}{}{}{
		\begin{itemize}
			\item Diseñé e implementé algoritmos de resolución numérica en Python para resolver ecuaciones diferenciales no lineales aplicadas a modelos físicos.
			\item Evalué la convergencia teórica y estabilidad computacional de métodos de descomposición matemática frente a esquemas tradicionales.
			\item \href{https://github.com/Maticus-7/lane-emden}{Haga click aquí para ver el repositorio en GitHub}
	\end{itemize}}
	
	
	\section{Actividades extracurriculares}
	\cvitem{}{Co-director del área de Proyectos e Investigación del IEEE Computational Intelligence Society – UNMSM (Gestión de proyectos tecnológicos y talleres de IA/ML).}
	\cvitem{}{Research Trainee en Bildheit (Investigación aplicada y modelación matemática).}
	\cvitem{}{Exmiembro del equipo universitario de judo – Universidad Nacional Mayor de San Marcos.}
	\cvitem{}{Exmiembro de la agrupación cultural y artística de salsa SALSUNI – San Marcos.}
	
	\section{Premios y Distinciones}
	\cvitem{}{BECA PERMANENCIA 2025 -- PRONABEC (Beca otorgada por excelencia académica a nivel nacional).}
	\cvitem{}{Reconocimiento de mi facultad por mi EXCELENCIA ACADÉMICA (Representando a mi universidad en eventos en el extranjero).}  
	
	\section{Eventos académicos}
	\cventry{Febrero de 2026}{UNT -- Trujillo}{Summer School and Workshop on Optimization and Operator Learning}{}{}{Puse a prueba mis conocimientos en teoría de optimización matemática y programación avanzada, destacando en el desarrollo de la pre-escuela.}
	\cventry{Diciembre de 2025}{Huaraz, Perú}{VII EICM}{}{}{Ponente del proyecto de investigación enfocado en el Modelamiento Numérico de Cuerpos Estelares.}
	\cventry{Octubre de 2025}{Concepción, Chile}{E'(UBB)=2025}{}{}{Participación internacional en ponencias y discusiones de matemáticas aplicadas y análisis numérico.}
	
	\section{Idiomas}
	\cvitem{Inglés}{Intermedio}
	\cvitem{Portugués}{Intermedio}
	
	\section{Certificados} 
	\cvitem{}{ \href{https://maticus-7.github.io/Certificates/}{Los certificados y logros académicos los podrá ver en mi portafolio web haciendo click aquí.}}
	
\end{document}